%%%%%%%%%%%%%%%%%%%%%%%%%%%%%%%%%%%%%%%%%%%%%%%%%%%%%%%%%%%%%%%%%%%%%%%%%%%%%%%%
%2345678901234567890123456789012345678901234567890123456789012345678901234567890
%        1         2         3         4         5         6         7         8

\documentclass[letterpaper, 10 pt, conference]{ieeeconf}  % Comment this line out if you need a4paper

%\documentclass[a4paper, 10pt, conference]{ieeeconf}      % Use this line for a4 paper

\IEEEoverridecommandlockouts                              % This command is only needed if 
                                                          % you want to use the \thanks command

\overrideIEEEmargins                                      % Needed to meet printer requirements.

%In case you encounter the following error:
%Error 1010 The PDF file may be corrupt (unable to open PDF file) OR
%Error 1000 An error occurred while parsing a contents stream. Unable to analyze the PDF file.
%This is a known problem with pdfLaTeX conversion filter. The file cannot be opened with acrobat reader
%Please use one of the alternatives below to circumvent this error by uncommenting one or the other
%\pdfobjcompresslevel=0
%\pdfminorversion=4

% See the \addtolength command later in the file to balance the column lengths
% on the last page of the document

% The following packages can be found on http:\\www.ctan.org
%\usepackage{graphics} % for pdf, bitmapped graphics files
%\usepackage{epsfig} % for postscript graphics files
%\usepackage{mathptmx} % assumes new font selection scheme installed
%\usepackage{times} % assumes new font selection scheme installed
%\usepackage{amsmath} % assumes amsmath package installed
\usepackage{amssymb}  % assumes amsmath package installed

\usepackage{bibentry}
% END REMOVE bibentry
\usepackage{multicol}
%\usepackage[bookmarks=true]{hyperref}
\usepackage{amsmath}
\usepackage{graphicx}
\usepackage{subfigure}
%\usepackage{amsthm} 
\usepackage{adjustbox}

 \newtheorem{definition}{Definition}

\title{\LARGE \bf
A Pipelined Framework for Multi-Robot Underwater 3D Inspection with Depth-Limited Pyramid Visibility
}


\author{Kryštof Pšenička$^{1}$  \and Stefan Edelkamp$^{1}$ \and Xuzhe Dang$^{2,1}$%
\thanks{$^{1}$% 
Faculty of Mathematics and Physics, 
Charles University, Prague, Czech Republic 
%{\tt\small edelkamp@mff.cuni.cz}
}
\thanks{$^{2}$ 
Faculty of Electrical Engineering,
Czech Technical University, Prague, Czech Republic 
}%
}

\begin{document}



\maketitle
\thispagestyle{empty}
\pagestyle{empty}


\begin{abstract}
We present a complete pipelined framework for multi-robot underwater inspection of 3D objects under depth-limited pyramid visibility. The pipeline consists of five stages: (1)~candidate viewpoint sampling around the target mesh, (2)~visibility computation via BVH-accelerated raycasting with pyramid-shaped view frustums, (3)~viewpoint selection formulated as a set cover problem and solved by a Greedy+Kernel optimizer, (4)~construction of a voxelized occupancy grid from the environment mesh and computation of a collision-free distance matrix using GPU-accelerated Dijkstra (cuGraph), and (5)~multi-robot route planning via a capacitated Vehicle Routing Problem solved with NVIDIA cuOpt, followed by cubic-spline trajectory execution with inter-robot collision avoidance. We evaluate the framework on the Duke of Lancaster shipwreck model in OceanSim/Isaac~Sim. A comparative analysis of visibility methods shows that raycasting is faster and more accurate than $\epsilon$-visibility, and among five set-cover optimizers the Greedy+Kernel method dominates the Pareto front of computation time versus viewpoint count.
\end{abstract}


%\IEEEpeerreviewmaketitle

\section{Introduction}

Inspection robots are automated devices equipped with advanced sensors and cameras designed to carry out inspections in environments that are hazardous, difficult to access, or require high precision. These robots are increasingly in demand because of their ability to perform tasks that are beyond the reach of human inspectors, often under harsh or dangerous conditions. The need for inspection robots has increased in recent years, driven by advances in robotics and the increasing complexity of industrial operations.

Some key sectors where 3D inspection robots are making a significant impact are: manufacturing, where robots are used for quality control, inspecting machinery, and monitoring production lines to identify defects and ensure that products meet stringent quality standards.
In the oil and gas sector, inspection robots are essential for monitoring pipelines, detecting leaks, and inspecting offshore rigs to prevent environmental hazards and reduce operational risks.

In our work, we consider underwater inspection missions with a fleet of multiple robots having limited visibility, mainly pyramid visibility and depth-limited range. We extend~\cite{pyramidvis} with depth-limited vision and multi-robot exploration, and~\cite{3dinspect} with limited sensors, GPU-accelerated vehicle routing, and a complete end-to-end pipeline.

Our main contributions are:
\begin{enumerate}
\item A complete pipelined framework for multi-robot 3D inspection under depth-limited pyramid visibility, covering every stage from viewpoint generation to collision-free route execution.
\item A comparative evaluation of visibility computation methods (raycasting vs.\ $\epsilon$-visibility) and five viewpoint selection optimizers, identifying Greedy+Kernel as Pareto-optimal.
\item A GPU-accelerated multi-robot routing module that constructs a voxelized occupancy grid, computes collision-free distances via cuGraph Dijkstra, and solves a capacitated Vehicle Routing Problem with NVIDIA cuOpt, producing balanced tours for each robot.
\end{enumerate}


\section{Related Work}

\subsection{3D Inspection and Visibility}
The 3D guarding problem asks for a minimum set of viewpoints from which the entire surface of a mesh is visible. Yu et al.~\cite{pyramidvis} introduced pyramid visibility constraints for autonomous robots, while star decomposition~\cite{Yu} and approximate star-shaped decomposition~\cite{PBG07_cr} provide theoretical foundations for visibility computation. Edelkamp et al.~\cite{3dinspect} formulated surface inspection as a hitting-set problem combined with multi-goal motion planning, and~\cite{watchman} studied watchman routes for single-robot inspection.

\subsection{Underwater Robot Planning}
Autonomous Underwater Vehicles (AUVs) pose unique planning challenges due to limited communication, constrained energy budgets, and degraded perception. McMahon et al.~\cite{inspectunderwater} addressed simultaneous survey and inspection with AUVs. Carreno et al.~\cite{Carreno_Pairet_Petillot_Petrick_2020} proposed a decentralized strategy for heterogeneous AUV missions via goal distribution and temporal planning. Hernandez et al.~\cite{Hernandez_Carreras_Ridao_2011} developed path planning for AUVs guided by homotopy classes. Pinto et al.~\cite{Pinto} described coordinated multi-robot ocean exploration. Song et al.~\cite{OceanSim} introduced OceanSim, a GPU-accelerated simulation framework for underwater robot perception that we use for validation.

\subsection{Vehicle Routing for Multi-Robot Systems}
The Vehicle Routing Problem (VRP) provides a natural formulation for assigning waypoints to multiple robots and sequencing their visits. GPU-accelerated VRP solvers such as NVIDIA cuOpt~\cite{cuopt} can handle fleet-level optimization with capacity constraints, enabling balanced route assignment. We combine cuOpt with cuGraph~\cite{cugraph} for computing collision-free distance matrices on voxelized 3D environments.

\section{Preliminaries}

Our approach follows a multi-stage pipeline. First, candidate inspection viewpoints are generated and their visibility of the target mesh is computed. The viewpoints are then filtered via set-cover optimization. Finally, the selected waypoints are assigned to robots and sequenced through a Vehicle Routing Problem, and collision-free trajectories are generated. Below we formalize the key concepts.

\subsection{Robot Model}

Each \emph{robot model} ${\cal R}_i$ is defined by its geometric shape ${\cal P}_i$, position $\mbox{\em Pos}(s_i) \in \mathbb{R}^3$, and orientation $\mbox{\em Orientation}(s_i)$. The function $\mbox{\em Place}({\cal P}_i, s_i)$ denotes the spatial footprint of robot~$i$ in state~$s_i$. Each robot has an initial position $s_{\text{init},i}$ (its depot) and a bounding radius~$r$ used for collision checking.

\subsection{Environment and Collision Checking}

The environment is defined by its bounding box $W$ and obstacles $O=\{O_1,\ldots,O_m\}$, where $O_j\subseteq W$. For path planning, the workspace is discretized into a 3D voxel occupancy grid at resolution~$\delta$. Each voxel is marked as occupied if it intersects the mesh or an obstacle. The grid is then \emph{inflated} by dilating occupied voxels with a spherical structuring element of radius $\lceil r/\delta \rceil + 1$ voxels, where $r$ is the robot bounding radius. This inflation allows robots to be treated as point agents during path planning while guaranteeing collision-free motion for the full robot geometry.

A composite state $(s_1,\ldots,s_n)$ is \emph{collision-free} if each robot is inside~$W$, no robot overlaps an obstacle, and no two robots overlap each other:
$\mbox{\em Place}({\cal P}_i,s_i) \cap \mbox{\em Place}({\cal P}_j,s_j) = \emptyset$ for all $i \neq j$.


\subsection{Visibility Regions and Pyramid Visibility}

Let $\mathcal{M} = \{M_1, \ldots, M_n\}$ denote the set of mesh regions to be inspected.
Let $\mbox{\em Vis}(M_i,p)$ denote the area of $M_i$ that is
visible from $p \in W$, i.e.,
$\mbox{\em Vis}(M_i,p) = \{q \mid q \in 
M_i \wedge \overline{pq} \cap \partial M_i = \emptyset\}.$ 
The visible area from a set of
points $(p_1 ,\ldots,p_l) \in W$ across all regions is
$\mbox{\em Vis}(\mathcal{M}, p_1,\ldots,p_l) =
 \bigcup_{M_i \in \mathcal{M}}
\bigcup_{j=1}^l \mbox{\em Vis}(M_i, p_j).$



Given a 3D region $M$, whose boundary is discretized by a 
mesh $\partial M=\{V,F\}$, where 
$V=\{v_1,v_2,\ldots,v_{|V|}\}$ are vertices and 
$F=\{f_1,f_2,\ldots,f_{|F|}\}$ are facets.  
%
%We can define the limited visible region by imposing specific pyramid visibility constraints. 
Given a point $p$, the spherical region originating at $p$ with radius $d$ can be approximated by a set of square pyramids $\{\Delta_p\}$. The height of the pyramid is $d$. Point $q$ is \emph{pyramid visible} to $\Delta^i_p$ if $\overline{pq}$ is located in the pyramid $\Delta^i_p$. 
%For simplification, we assume that the spherical radius $d$ and the edge-length of square of the pyramids are fixed; so that the pyramids on each skeleton node can be pre-computed. 
Each pyramid $\Delta^i_p$ can be represented by $(\vec{p},\vec{d}_i)$,where $\vec{p}$ is the apex, $\vec{d}_i$ is the vector from $\vec{p}$ to the center of the base. The guarding problem with pyramid visibility is to find the minimal set of pyramids to cover all the boundary vertices. 
%The time complexity for computing the limited visible region can be reduced with a scan-line. %~\cite{}.

\subsubsection{Theoretical Visibility Models}
To solve the guarding problem effectively, we must first establish a method for determining which parts of the mesh are visible from a candidate viewpoint. Two prominent approaches derived from computational geometry and the star-decomposition literature are $\epsilon$-visibility and sweeping visibility.

\begin{itemize}
    \item \textbf{$\epsilon$-Visibility}: Also known as approximate star-shaped decomposition, this method is particularly robust for processing dense Point Set Data (PSD)~\cite{PBG07_cr}. A 3D model is considered star-shaped if there exists a kernel (a set of points) from which the entire model is visible. $\epsilon$-visibility relaxes this strict requirement; for a viewpoint $p$, the visible region $\text{Vis}_\epsilon(p)$ includes all mesh points $q$ where the line segment $\overline{pq}$ does not intersect the mesh boundary, allowing for a small tolerance $\epsilon$. This tolerance makes the algorithm resilient to noise in the mesh data, as it depends primarily on point positions and outward normals rather than strict connectivity.
    
    \item \textbf{Sweeping Visibility}: This approach provides an exact solution~\cite{Yu}. The algorithm employs a rotational plane sweep or spherical sweep around the viewpoint $p$. As the plane sweeps across the angular domain, the algorithm maintains a status structure of the mesh faces intersecting the plane, allowing precise determination of occlusions.
\end{itemize}

\subsection{Ray Casting and Bounding Volume Hierarchy}

Instead, we work on \emph{ray casting} and \emph{ray tracing}.
A \emph{ray} is a half-line in 3D defined by an origin $\vec{p} = (x_0, y_0, z_0)$ and a direction vector $\vec{d} =(d_x, d_y, d_z)$. Points on the ray are parameterized as $\vec{p} + t\vec{d}$ for $t \ge 0$.
If we normalize $\hat{d} = \vec{d}/\|\vec{d}\|_2$, then $t$ equals the Euclidean distance along the ray.
The \emph{ray casting problem} asks, given a set of objects, to preprocess them into a data structure so that for each query ray the first object hit can be found quickly. 

A \emph{bounding volume hierarchy} (BVH) is a tree-like data structure used in computer graphics to accelerate intersection tests. We organize objects by grouping them into a hierarchy of bounding volumes, starting with a single root volume that encloses all objects in the scene. When a ray or collision query is made, the structure allows the system to quickly determine if the query intersects any objects by first checking the bounding volumes. 

\subsection{Waypoint Filtering via Capacitated Set Cover}

Once the visibility information for all sampled viewpoints is computed, we select a minimum-sized set of viewpoints that covers all boundary vertices. 
For this optimization, we formulate \emph{waypoint filtering} as a capacitated set cover problem. It is well-known that set covers are polynomially equivalent to hitting sets in bipartite graphs.
In the non-capacitated $\mbox {\em Set  Cover}$ instance, we are given a universe $X$ of $n$ elements and a collection $V$ of $m$ subsets of $X$, each subset with an associated cost; the goal is
to choose the collection of sets $V' \subseteq V$ of the least total cost, such that each element $e \in X$ is
contained in at least one set $V \in V'$. 
%
In an \emph{Edge-weighted Capacitated Set Cover} ECSC$(A, d, v, c)$ we optimize $$\min \sum_{i=1} c_i x_i,\mbox{ s.t. } 
\sum_{i=1} A_{i,j} v_{i,j} x_i \ge d_j, \mbox{for all } 1 \le j \le n$$ 
and 
$x_i \in 
\{0, 1\} \mbox{ for all } i \in \{1,\ldots,n\}$. In our application, $d_j$ is the visibility demand for mesh node~$j$ (e.g., 95\%), $v_{i,j}$ is the visibility score of mesh point~$j$ as seen from viewpoint~$i$, $x_i$ are the indicator variables selecting viewpoint~$i$, $c_i$ is the associated cost (e.g., travel distance to reach the position), and $A_{i,j}$ is the binary selection matrix.

\section{Problem Formulation}

In the simplest case, the quality of a set of inspection points 
$p_1 ,\ldots,p_l \in W$ is defined as the
fraction of the mesh area in $\mathcal{M}$ that is visible:
$$
\mbox{\em Quality}(\mathcal{M}, p_1,\ldots,p_l) = \frac{\mbox{\em Area}(\mbox{\em Vis}(\mathcal{M}, p_1 ,\ldots,p_l)) 
}
{\mbox{\em Area}(\bigcup_{M_i \in \mathcal{M}} M_i)}$$
In more advanced settings, we assign \emph{visibility scores} to each 
mesh unit with respect to every sampled robot
position and orientation. 
\emph{Quality} is then defined by the normalized 
solution of the ECSC:
$$
\mbox{\em Quality}(\mathcal{M}, p_1,\ldots,p_l) = \mbox{\em VisScore}(\mathcal{M}, p_1 ,\ldots,p_l).$$
The inspection problem for multiple robots can now be stated as follows.

\begin{definition}
Given a desired inspection quality $0 < \alpha \le 1$, an environment $W$, obstacles $O = \{O_1,\ldots, O_m\}$, mesh regions
$\mathcal{M} = \{M_1,\ldots, M_n \}$, and $k$ robots with starting positions $s_{\text{init},i} \in W$ for $i\in \{1,\ldots,k\}$, compute
\begin{itemize}
\item a set of inspection viewpoints 
$P = \{p_1 ,\ldots,p_l\} \subseteq W$ such that
$\mbox{\em Quality}(\mathcal{M}, p_1 ,\ldots,p_l) \ge \alpha$,
\item cost-minimizing open tours $\pi_i = (s_{\text{init},i}, p^i_1 ,\ldots,p^i_{l_i})$, $i\in \{1,\ldots,k\}$, starting at each robot's depot,
\item collision-free trajectories $\rho_i$:
$\{1,\ldots,T\} \rightarrow W$ for all $i\in \{1,\ldots,k\}$, which visit the waypoints
in the order defined by $\pi_i$ such that all trajectories are mutually collision-free.
\end{itemize}
\end{definition}

\section{Pipeline}

Our framework solves the multi-robot inspection problem through a five-stage pipeline. In the following, we describe each stage.

\subsection{Stage 1: Viewpoint Sampling}

Candidate inspection viewpoints are generated around the target mesh. For each candidate position~$p$ and orientation, a pyramid-shaped view frustum with depth limit~$d$ is constructed. The union of mesh vertices falling inside the frustum and not occluded by other mesh faces constitutes the visibility set of that viewpoint. We use a $kd$-tree for fast frustum culling and BVH-accelerated raycasting for occlusion testing. Fig.~\ref{fig:visibilitymult} shows that this implementation computes visibility for 200 viewpoints on a point cloud with 1\,000\,000 points in approximately 70\,s on a 2.2\,GHz Intel Core i5.

\subsection{Stage 2: Viewpoint Selection}

Given the visibility sets from Stage~1, we solve the ECSC formulated in Section~3 to select a minimum-cost subset of viewpoints achieving the desired coverage~$\alpha$. As shown in our experimental evaluation (Section~6), the Greedy+Kernel optimizer provides the best trade-off between solution quality and computation time.

\subsection{Stage 3: Occupancy Grid Construction}

The environment mesh (including the inspection target and static obstacles such as rocks and coral) is voxelized. The resulting occupancy grid is then \emph{inflated} by dilating occupied voxels with a spherical structuring element of radius $\lceil r/\delta \rceil + 1$ voxels, where $r$ is the robot bounding radius. After inflation, robots can be treated as point agents during path planning.

\subsection{Stage 4: Multi-Robot Route Planning via VRP}

\subsubsection{Distance Matrix}
A 26-connected undirected graph is constructed from the free voxels of the inflated occupancy grid, where edge weights equal the Euclidean distance between neighboring voxel centers scaled by the voxel resolution. All-pairs shortest-path distances between the $N$ selected waypoints (plus per-robot depot positions) are computed on the GPU using NVIDIA cuGraph Dijkstra~\cite{cugraph}.

\subsubsection{Capacitated VRP}
The resulting $N \times N$ cost matrix is passed to NVIDIA cuOpt~\cite{cuopt}, which solves a capacitated Vehicle Routing Problem. Each non-depot waypoint carries a unit demand, and each vehicle is assigned capacity $\lceil N_{\text{wp}} / k \rceil + 1$, where $N_{\text{wp}}$ is the number of inspection waypoints and $k$ is the fleet size. This capacity constraint ensures balanced routes. Per-vehicle depot indices support heterogeneous start locations.

\subsection{Stage 5: Trajectory Execution and Collision Avoidance}

The shortest voxel-grid paths underlying the distance matrix are extracted and converted to world-frame waypoint sequences (A* path cache). Each route leg is interpolated at constant cruise speed with arc-length-parameterized \emph{cubic spline} fitting for legs with four or more keypoints, and linear interpolation otherwise. Inter-robot collisions are resolved via a priority-based scheme: at each time step, axis-aligned bounding boxes (AABBs) of all robots are checked for overlap; when a conflict is detected, the lower-priority robot (higher index) is held in place for the duration of the overlap window. This process iterates until no collisions remain.

The decision to handle collisions \emph{a posteriori} rather than planning jointly collision-free multi-robot trajectories is justified by two observations: first, the robot density is low relative to the environment volume (we use at most 6 AUVs); second, the VRP solution naturally assigns each robot a spatially clustered subset of waypoints, so their paths rarely intersect. This reactive scheme therefore avoids the combinatorial cost of coupled multi-robot motion planning while still guaranteeing collision-free execution in practice.


\begin{figure}[t]
    \centering
    \includegraphics[width=0.8\linewidth]{visibility.png}
    \caption{Combined visibility of a sunken ship Duke of Lancaster with about 1 million mesh points from 200 random viewpoints, green areas are visible while red areas are not.}
    \label{fig:visibilitymult}
\end{figure}

\section{Experimental Evaluation}

We apply our pipelined framework to underwater
inspection tasks using scenarios from OceanSim~\cite{OceanSim} based on NVIDIA IsaacSim. For the mesh to be inspected, we generate candidate viewpoints and evaluate pyramidal visibility for each of them using raycasting with BVH acceleration in Open3D, followed by solving an edge-weighted set cover problem.

\subsection{Rationale for Excluding Sweeping Visibility}

While the sweeping visibility algorithm provides mathematically exact solutions for 3D guarding, its computational complexity makes it ill-suited for the specific constraints of our setting. The primary bottleneck is the nature of the \textbf{Pyramid Visibility} constraint.

In our inspection model, visibility is restricted to a depth-limited view frustum. Algorithms such as \textbf{Raycasting} and \textbf{$\epsilon$-Visibility} benefit significantly from this restriction through \textit{geometric culling}. By using spatial partitioning (e.g., BVH or kd-trees), these methods can instantly discard the vast majority of mesh vertices that fall outside the view frustum, reducing the complexity to significantly fewer points.

In contrast, the \textbf{Sweeping Visibility} algorithm operates by constructing a visibility map based on the global topology of the mesh to determine event points and critical triangles. Even if the final result is clipped to the pyramid view, the internal construction of the sweeping plane or spherical map necessitates traversing or processing the full set of triangle vertices to resolve occlusions correctly. Consequently, the sweeping algorithm does not achieve the same order-of-magnitude speedup from the pyramid constraint as the other methods. Given this lack of scalability with restricted fields of view, we excluded the sweeping algorithm from the following comparative analysis.

\subsection{Evaluation of Visibility Methods}

We evaluated the performance of $\epsilon$-visibility against Raycasting to get visibility of the mesh with 200,000 points from 2,000 inspection viewpoints on the ``Duke of Lancaster'' ship model. The objective was to determine which method is more efficient across multiple coverage targets. 

\begin{figure}[t]
    \centering
    \includegraphics[width=0.75\linewidth]{Duke Of Lancaster_Coverage_Bar.png}
    \caption{Coverage for varying coverage targets.}
    \label{fig:coverage_comparison}
\end{figure}

\begin{figure}[t]     
    \centering
    \includegraphics[width=0.75\linewidth]{Duke Of Lancaster_Viewpoint_Count_Bar.png}

    \caption{Number of viewpoints wrt.\ coverage targets.}
    \label{fig:viewpoint_count}
\end{figure}

\subsubsection{Viewpoint Efficiency}
The actual coverage achieved by the approximate $\epsilon$-visibility algorithm is computed using the raycasting implementation. As illustrated in Fig.~\ref{fig:coverage_comparison}, the $\epsilon$-visibility implementation does not exactly meet the coverage targets because it is slightly inaccurate (for some viewpoints it overestimates visibility and for others it underestimates it). In addition, Fig.~\ref{fig:viewpoint_count} reveals a critical difference in efficiency: \textbf{Raycasting consistently required fewer viewpoints} to achieve the same (actually higher) coverage. This suggests that the approximate nature of $\epsilon$-visibility forces the viewpoint selection algorithm to choose multiple redundant viewpoints that cover points that are "unknowingly" covered by other viewpoints.

% \begin{figure}[t]
%     \centering
%     \begin{adjustbox}{max width=0.55\linewidth}
%         \includegraphics[width=0.55\textwidth]{Duke Of Lancaster_Visibility_Time_Bar.png}
%     \end{adjustbox}
%     \caption{Computational time analysis of visibility computation for 2000 viewpoints on mesh with 200 000 target points.}
%         \label{fig:visibility_time}
% \end{figure}

\subsubsection{Computational Performance}
The most significant disparity was observed in the computation time. The raycasting method took 79.6s and the $\epsilon$-visibility method took 153.3s. This performance gap is attributed to the overhead of maintaining the approximate star-shaped decomposition structures versus the speed of hardware-optimized ray-triangle intersection tests.

This analysis clearly demonstrates that Raycasting offers a superior trade-off between speed and solution quality for this class of inspection problems.

\subsection{Evaluation of Viewpoint Selection Methods}

After computing visibility using raycasting, we select an optimal set of viewpoints to achieve the desired coverage. This is formulated as a set cover problem. We evaluated five optimization algorithms: \textit{Greedy} (iteratively selects viewpoints covering the most uncovered points), \textit{Greedy+Kernel} (greedy with random local neighborhood search to refine each selection), \textit{NRPA}~\cite{nrpa} (Monte Carlo tree search with learned rollout policies), \textit{Beam Search} (NRPA variant maintaining multiple partial solutions), and \textit{Nested Monte Carlo} (recursive nested search).
%
Experiments in the \emph{Duke of Lancaster model} used 100,000 target points, 1,000 candidate viewpoints, and 95\% target coverage.

\subsubsection{Results}

MCTS methods (NRPA, Beam) find valid solutions quickly (Fig.~\ref{fig:phase_a_time}), but require $4$-$7$ times more viewpoints than greedy approaches (Fig.~\ref{fig:phase_a_viewpoints}). This occurs because MCTS optimizes for speed rather than solution size---random rollouts select viewpoints covering small increments, creating redundant solutions. Greedy methods make locally optimal choices that compound into near-optimal solutions.
%
Fig.~\ref{fig:optimizer_pareto} confirms Greedy+Kernel dominates the Pareto front.

\textbf{Best Choice}: For minimizing viewpoints, directly reducing mission duration and energy consumption, \textit{Greedy+Kernel} is preferred. It requires twice the time of pure Greedy (75s vs.\ 34s) but produces 23\% fewer viewpoints (58 vs.\ 75). MCTS methods are unsuitable for this problem structure.

\begin{figure}[t]
    \centering
        \includegraphics[width=0.8\linewidth]{optimizer_comparison_graphs/pareto_fronts.png}
    \caption{Pareto fronts showing the trade-off between number of viewpoints and coverage for each method. Greedy+Kernel dominates the Pareto front.}
    \label{fig:optimizer_pareto}
\end{figure}

\begin{figure}[t]
    \centering
    \subfigure[]{
            \includegraphics[width=0.45\linewidth]{optimizer_comparison_graphs/phase_a_time_to_coverage.png}
        \label{fig:phase_a_time}
    }%
    \subfigure[]{
        \includegraphics[width=0.50\linewidth]{optimizer_comparison_graphs/phase_a_viewpoints_at_coverage.png}
        \label{fig:phase_a_viewpoints}
    }
    \caption{Time and viewpoint count for the first valid solution.}
    
\end{figure}

\subsubsection{Best choice}

For inspection applications where \textit{minimizing the number of viewpoints} is critical, which directly translates to shorter mission duration, reduced energy consumption, and lower operational costs, the \textit{ Greedy + Kernel} method is clearly preferred. Although it requires approximately twice the optimization time of pure Greedy ($75$~s vs.\ $34$~s for 95\% coverage), it produces solutions with 23\% fewer viewpoints ($58$ vs.\ $75$).

The MCTS-based methods, while theoretically appealing for their anytime properties and ability to escape local optima, are not well-suited for this problem structure. The hitting set problem with a given target coverage does not benefit from the exploratory nature of Monte Carlo sampling; instead, the greedy heuristic's locally optimal choices compound into globally near-optimal solutions.

\subsection{Multi-Robot Routing}

We evaluate the VRP-based routing module by sweeping fleet sizes $k \in \{1,\ldots,6\}$ and waypoint counts $N \in \{10, 20, 30, 50, 75, 100\}$, averaging over three random seeds per configuration. All 108 runs used the cuOpt solver on the Duke of Lancaster occupancy grid at $\delta = 0.10$\,m resolution.

\subsubsection{Makespan Reduction}
Fig.~\ref{fig:vrp_makespan} shows the longest single-robot route (makespan) as a function of fleet size. For 100 inspection waypoints, a single robot must traverse $367.6$\,m; with $k{=}4$ robots the makespan drops to $119.7$\,m --- a $3.07\times$ speedup (Fig.~\ref{fig:vrp_speedup}). Scaling beyond four robots yields diminishing returns: $k{=}6$ achieves only $3.09\times$ speedup with high variability.

\begin{figure}[t]
    \centering
    \subfigure[Makespan vs.\ fleet size]{%
        \label{fig:vrp_makespan}%
        \includegraphics[width=0.48\linewidth]{fig_makespan_vs_fleet.png}%
    }%
    \subfigure[Speedup vs.\ fleet size]{%
        \label{fig:vrp_speedup}%
        \includegraphics[width=0.48\linewidth]{fig_speedup_vs_fleet.png}%
    }
    \caption{Makespan and speedup relative to a single robot. Dashed line: ideal linear speedup. Increasing the fleet to 4 robots yields a $3.07\times$ makespan reduction on 100 waypoints; further scaling provides marginal improvement.}
    \label{fig:vrp_fleet_scaling}
\end{figure}

\subsubsection{Computation Time}
Fig.~\ref{fig:vrp_timing} shows the per-stage time breakdown for 50 waypoints. The distance-matrix computation (${\sim}135$\,s) dominates at all fleet sizes, with the cuOpt solver adding up to $62$\,s for multi-robot instances. The A* path cache scales with fleet size as more route segments are computed. Overall, the pipeline completes 100-waypoint, 6-robot instances in under $484$\,s (${\sim}8$\,min), with the single-robot baseline at $280$\,s.

Fig.~\ref{fig:vrp_scalability} plots total time against the number of waypoints. Scaling is sub-linear: a $10\times$ increase from 10 to 100 waypoints roughly doubles the distance-matrix time ($93 \to 190$\,s).

\begin{figure}[t]
    \centering
    \includegraphics[width=0.75\linewidth]{fig_timing_breakdown.png}
    \caption{Computation time breakdown by pipeline stage (50 waypoints). The GPU distance matrix dominates for small fleets; the VRP solver and A* path cache grow with fleet size.}
    \label{fig:vrp_timing}
\end{figure}

\begin{figure}[t]
    \centering
    \includegraphics[width=0.75\linewidth]{fig_scalability_waypoints.png}
    \caption{Total pipeline time vs.\ waypoint count, by fleet size. Scaling is sub-linear in the number of waypoints due to fixed grid-construction overhead.}
    \label{fig:vrp_scalability}
\end{figure}


% \section{Simulation}

% For multi-robot collision free path planning between the waypoints, we adapted cuRobo for real-time obstacle-free movements of industrial arms. For each UAV, instead of an arm, we defined a robot model that includes invisible and non-colliding three prismatic joints along the axis together with rotation (6 DOF), with the underwater vehicle as end effector. We also have put a prismatic camera in front of the AUV to enforce heading as each waypoint is stored as a 7-tuple. 

% We work on centralized single-robot (see Fig.~\ref{fig:underwatersimone}) and multi-robot path planning  (see Fig.~\ref{fig:underwatersimmult}).
% Robot-robot collision is a problem in space-time. We are approximated robots by bounding boxes and for a fixed ordering of the robots we block their traces for the next robot to bypass (see Fig.~\ref{fig:underwatersimmult}). Refinements of this coarse strategy are possible, but as VRP solvers tend to cluster the waypoints into different regions where robot-robot collisions are rare, we did not elaborate on computing the exact solution. 

  

% \begin{figure}[t]
%     \centering
%     \includegraphics[width=0.48\linewidth]{ underwater4.png}
%     \includegraphics[width=0.48\linewidth]{ underwater2.png}
%     \includegraphics[width=0.48\linewidth]{ underwater3.png}
%     \includegraphics[width=0.48\linewidth]{ underwatersim.png}

%     \caption[screenshots of simulation]{Connecting waypoints for simulation of submarine movements avoiding obstacles.}
%     \label{fig:underwatersimone}
% \end{figure}


% \begin{figure}[t]
%     \centering
%     \includegraphics[width=0.68\linewidth]{underwatermultirob.png}
    
%     \caption {Connecting 50 colored waypoints in simulation of five submarines avoiding robot-to-robot and robot-to-obstacles collisions.}
%     \label{fig:underwatersimmult}
% \end{figure}

% \begin{figure}[t]
%     \centering
%     \includegraphics[width=0.3\linewidth]{collision1.png}
% \includegraphics[width=0.3\linewidth]{collision2.png}
% \includegraphics[width=0.3\linewidth]{collision3.png}
%     \caption{Temporal-spatial  incremental robot-to-robot collision detection: (left) path for first red robot and bounding box (squares are start and target position and ball illustrates orientation), 
%     (middle) conflict-free path for second blue robot with first one with trace of first one (shown in yellow) is blocking the way and has to be bypassed. (right) path for first robot traveling back with second robot stationary at its target position. }
%     \label{fig:collision}
% \end{figure}

\section{Conclusion and Future Work}

We presented a complete pipelined framework for multi-robot underwater 3D inspection under depth-limited pyramid visibility. The pipeline spans five stages: viewpoint sampling, visibility evaluation using BVH-accelerated raycasting, viewpoint selection via capacitated set cover optimization, voxelized occupancy grid construction, GPU-accelerated multi-robot route planning using cuGraph and cuOpt, and collision-free trajectory execution with inter-robot deconfliction.

Our experimental evaluation on the Duke of Lancaster shipwreck model yielded three key findings. First, raycasting provides a superior trade-off between speed and accuracy compared to $\epsilon$-visibility for pyramid-constrained visibility computation. Second, among five set-cover optimizers, the Greedy+Kernel method dominates the Pareto front, producing 23\% fewer viewpoints than pure greedy selection at the cost of roughly doubled optimization time. Third, the VRP-based routing module achieves a $3.07\times$ makespan reduction with four robots on 100-waypoint instances, with the full pipeline completing in under 13 minutes. For the multi-robot simulation of the results we refer to the video (for reviewers: to be found in the supplementary material, later on in case of acceptance: web-address)

In future work, we plan to incorporate \emph{location-routing}, where the cost of reaching a viewpoint influences its selection during set cover optimization. We also aim to evaluate the framework on larger and more complex environments, to integrate energy-aware planning that accounts for ocean currents and varying terrain difficulty, and to improve the visibility computation by extending HPRO~\cite{HRPO}, a differentiable visibility operator that maps the problem to convex-hull computation.

%\nocite{*}

\bibliographystyle{plain}
\bibliography{main}


\end{document}



\begin{thebibliography}{99}
\bibitem{pyramidvis}
Wuyi Yu, Maoqing Li, Xin Li: \emph{Optimizing Pyramid Visibility Coverage for Autonomous robots in 3D Environment}. Control. Intell. Syst. 42(1) (2014)
\bibitem{inspectunderwater}
James McMahon, Riley Parker, Philip D. Baldoni, Stuart Anstee, Erion Plaku: \emph{Simultaneous Survey and Inspection with Autonomous Underwater Vehicles}. IROS 2023: 1-7
\bibitem{dromosmr}
Duong Le, Erion Plaku: \emph{Multi-Robot Motion Planning With Dynamics via Coordinated Sampling-Based Expansion Guided by Multi-Agent Search}. IEEE Robotics Autom. Lett. 4(2): 1868-1875 (2019)
\bibitem{dromostw}
Stefan Edelkamp, Morteza Lahijanian, Daniele Magazzeni, Erion Plaku:
\emph{Integrating Temporal Reasoning and Sampling-Based Motion Planning for Multigoal Problems With Dynamics and Time Windows}. IEEE Robotics Autom. Lett. 3(4): 3473-3480 (2018)
\bibitem{OceanSim}
Jingyu Song, Haoyu Ma, Onur Bagoren, Advaith V. Sethuraman, Yiting Zhang and Katherine A. Skinner. \emph{OceanSim: A GPU-Accelerated Underwater Robot Perception Simulation Framework}, https://arxiv.org/abs/2503.01074 (2025)
\bibitem{nrpa} 
Christopher D. Rosin: \emph{Nested Rollout Policy Adaptation for Monte Carlo Tree Search}. IJCAI 2011: 649-654
\bibitem{watchman} 
Stefan Edelkamp, Zhuowei Yu: \emph{Watchman Routes for Robot Inspection}. TAROS (2) 2019: 179-190
\bibitem{3dinspect} 
Stefan Edelkamp, Baris Can Secim, Erion Plaku: \emph{Surface Inspection via Hitting Sets and Multi-goal Motion Planning}. TAROS 2017: 134-149
\bibitem{energymr} 
Jachym Herynek and Stefan Edelkamp, \emph{Multi-Robot Multi-Goal Mission Planning in Terrains of Varying Energy Consumption}, IROS: 12064-12069 (2024)
\bibitem{prm} 
Kavraki, L.E., P. Svestka, J.-C. Latombe, and M.H. Overmars. \emph{Probabilistic roadmaps for path planning in high-dimensional configuration spaces}, IEEE Transactions on Robotics and Automation. Vol. 12, No. 4, Aug 1996 pp. 566—580.
\bibitem{Computer Graphics Forum - 2011 - Yu - Computing 3D Shape Guarding and Star Decomposition}
W. Yu and X. Li: \emph{Computing 3D Shape Guarding and Star Decomposition}. Computer Graphics Forum 30(7) (2011)
\bibitem{PBG07_cr}
Jyh-Ming Lien: \emph{Approximate Star-Shaped Decomposition of Point Set Data}. Eurographics Symposium on Point-Based Graphics (2007)
\end{thebibliography}




For the initial assignments of goals to the robots we use the NRPA algorithm
as an anytime solution and can be controlled by its parameters: the number of iterations (the width of the recursion) and the level of the search (the depth of the recursion tree). As every rollout yields an assignment of goals to agents the algorithm is anytime, and warrants an goal agenda assignment after a polynomial time steps.  
The implementation of the recursive search procedure for shortest-path reduced map is provided in Fig.~\ref{fig:nrpa}.
\begin{figure}[t]
    \centering  
\begin{small}
\begin{tabbing} 
\quad\=\quad\=\quad\=\quad\=\quad\=\kill
{\sc Search}($level,iteration$)  \\
  \>$best.score \leftarrow \infty$; \\
  \>{\bf if} ($level = 0$)  \\
  \>\> $best.score \leftarrow$ {\sc Rollout}(); \\
  \>\> for ($j = 0..m+k-1$) $best.tour[j] \leftarrow tour[j]$;  \\
  \> {\bf else}  \\
  \>\> $policy[level] \leftarrow global$; \\
  \>\> {\bf for} ($i=0..iteration-1$)  \\
  \>\>\> $r \leftarrow {\sc Search}(level - 1,iteration$); \\
  \>\>\> {\bf if} ($r.score < best.score$)  \\
        \>\>\>\> $best.score \leftarrow r.score$; \\
        \>\>\>\> for ($j = 0..m+k-1$) $best.tour[j] \leftarrow r.tour[j]$;  \\
        \>\>\>\> {\sc Adapt}($best.tour,level$); \\
  \>\>$global \leftarrow policy[level]$;  \\
  \> {\bf return} $best$;  \\
\end{tabbing}
\end{small}
\caption{Main search procedure NRPA for computing initial assignments in multi-goal multi-agent search. The parameter level is counted down, while the iteration width is kept constant. The policy is kept in each level and refreshed when going down the recursion. }
\label{fig:nrpa}
\end{figure} 
\begin{figure}[t]
    \centering  
\begin{small}
\begin{tabbing}
\quad\=\quad\=\quad\=\quad\=\quad\=\kill
{\sc Adapt}(tour, level) \\
  \> {\bf for} ($j=1..m$) $visits[j] \leftarrow 1$; \\
  \> visits[0] $\leftarrow$ k-1; succs $\leftarrow$ 0; node $\leftarrow$ 0; \\
  \> {\bf for} ($j = 0..m+k-1$)  \\ 
    \>\> $succs \leftarrow 0$ \\
    \>\> {\bf for} ($i = 0..m+k-1$) \\
    \>\>\>  {\bf if} ($visits[i] \neq 0$) $moves[succs] \leftarrow i; succs \uparrow 1$; \\
    \>\> $\lambda \leftarrow (node = 0) ? 1 : 1/k$ \\   
    \>\> z $\leftarrow$ 0; \\
    \>\> $policy[level][node][tour[j]] \uparrow \lambda$ \\
    \>\> {\bf for} ($i=0..succs-1$) $z \uparrow e^{global[node][moves[i]]}$; \\
    \>\> {\bf for} ($i=0..succs-1$) \\
    \>\>\>  $policy[level][node][moves[i]] \downarrow \lambda \cdot
        e^{(global[node][moves[i]])/z}$; \\
    \>\> $node \leftarrow tour[j]$; \\
    \>\> $visits[node] \downarrow 1$;  \\
\end{tabbing}
\end{small}
\caption{Policy adaptation in NRPA changing the policy in a given search level based on an improved tour. We use counters for the number of visits at each node, $\lambda$ is the learning parameter, $\uparrow$ is an increase, $\downarrow$ a decrease in value. 
In each iteration, the successors are generated and stored in an array, before being processed to tune the values in the policy matrix (increase for a good selection, decrease for a bad one).
}
\label{fig:adapt}
\end{figure}
The randomized simulation procedure or rollout
that is executed on each leaf of the motion tree is implemented in Fig.~\ref{fig:rollout}. It computes the sum-of-costs of the individual routes taken and returns it as an evaluation. The procedure that adapts the policy and steers the random simulation is shown in Fig.~\ref{fig:adapt}. The core implementation trick is to generate one tour for all of the agents in common and split it at the next obtained starting locations. When generating the combined tour, the distance from the last location of the current agent to the next is reset, and the makespan, too.
%    if (successors == 0) {
%      for(int i = 0; i < GOALS; i++)
%        if(check(i))
%           moves[successors++] = i; }

\begin{figure}[t]
    \centering 
\begin{small}
\begin{tabbing}
\quad\=\quad\=\quad\=\quad\=\quad\=\quad\=\quad\=\kill 
{\sc Rollout}() \\
\>  {\bf for} ($j=0..k-1$)    $visits[j] \leftarrow  1$; \\
\>  $visits[0] \leftarrow  k-1;  
tour[0] \leftarrow  0$; 
$size \leftarrow  1$;  $node \leftarrow  0$, $prev \leftarrow  0$; \\
\>  $ms \leftarrow 0$, $cost \leftarrow 0$; 
$span \leftarrow 0$; \\
\>  {\bf while} ($size < m+k-1$) \\
\>    \>$succs \leftarrow 0$; \\
\>    \>{\bf for} ($i = 0..m+k-1$)  \\
\>    \>  \>{\bf if} ($visits[i] \neq 0$) $moves[succs] = i; succs \uparrow 1$ \\
%\>    \>    \>\>{\bf for} (j = 0..k-1) \\
%\>    \>    \>  \>\>{\bf if} (i != j) \\
%\>    \>\>        \>\>\>{\bf if} (visited[j] = 0) \\
%\>    \>\>         \>\>\> \>succs--; break; \\ 
    \>\>{\bf for} ($i=0..succs-1$)  $value[i] \leftarrow e^{global[node][moves[i]]}$; \\ 
    \>\>$i \leftarrow$ {\sc Roulette}($value$); \\  
    \>\>$prev \leftarrow node; node \leftarrow moves[i]$; \\
    \>\> $tour[size] \leftarrow node; size \uparrow 1; visits[node] \downarrow 1$; \\
    \>\> $cost \uparrow d[prev][node]; span \uparrow d[prev][node]$; \\
    \>\>{\bf for} ($r = 0..k-1$) \\
      \>\>\>{\bf if} ($start[r] = node$) \\
        \>\>\>\>$cost \downarrow d[prev][node]$; \\
        \>\>\>\>$ms \leftarrow \max(span,ms); span \leftarrow 0$; \\
        \>\>\>\> {\bf break}; \\
\>  $ms \leftarrow \max(span,ms)$ \\ 
\>  {\bf return} $ms + cost/k$; \\
\end{tabbing}
\end{small}
\caption{Polynomial-time rollout procedure in NRPA. We use counters for the number of visits at each node, 
$d$ is the precomputed shortest path distance between two nodes,
$\lambda$ is the learning parameter (usually 1), $\uparrow$ is an increase, $\downarrow$ a decrease in value. In each iteration, the successors are generated and stored in an array, before being processed to tune the values in the policy matrix. In each iteration, the successors are generated and stored in an array, before being processed. Function Roulette determines the valid successor according to his share in the fitness value. A combination between makespan and cost is returned as objective.}
\label{fig:rollout}
\end{figure} 



\section{Experiments}
